% !TeX root = ../sections/02_exercises.tex

\begin{exercise}
	R-1. 環 $\mathbb{Z}[\sqrt{2}]$ は，$\mathbb{Z}$-加群とみると階数 $2$ の
	$\mathbb{Z}$-自由加群であることを示せ。
\end{exercise}

\begin{background}
	$\mathbb{Z}$-加群とは，アーベル群に整数倍が定まったもののことである。
	特に，$\mathbb{Z}$-自由加群であることを示すには，ある元の組が
	$\mathbb{Z}$ 上の基底になっていることを示せばよい。
	
	$\mathbb{Z}$-加群 $M$ に対して，元の組
	\[
	v_1,\dots,v_r\in M
	\]
	が基底であるとは，任意の $x\in M$ が一意的に
	\[
	x=a_1v_1+\cdots+a_rv_r
	\qquad
	(a_1,\dots,a_r\in\mathbb{Z})
	\]
	と表されることをいう。
	
	このとき $M$ は階数 $r$ の自由 $\mathbb{Z}$-加群である。
\end{background}

\begin{strategy}
	$\mathbb{Z}[\sqrt{2}]$ の元は，定義より
	\[
	a+b\sqrt{2}
	\qquad
	(a,b\in\mathbb{Z})
	\]
	の形で表される。
	したがって，まず $\{1,\sqrt{2}\}$ が
	$\mathbb{Z}[\sqrt{2}]$ を生成することを確認する。
	
	次に，一意性を示すために
	\[
	a+b\sqrt{2}=0
	\qquad
	(a,b\in\mathbb{Z})
	\]
	ならば $a=b=0$ であることを示す。
	ここで $\sqrt{2}$ が無理数であることを用いる。
	
	以上により，$\{1,\sqrt{2}\}$ が
	$\mathbb{Z}[\sqrt{2}]$ の $\mathbb{Z}$-基底であると分かる。
	したがって，$\mathbb{Z}[\sqrt{2}]$ は階数 $2$ の
	$\mathbb{Z}$-自由加群である。
\end{strategy}